\section{Discussion}
\label{sec:discussion}

This study tested whether a geometry-preserving generative method could produce
synthetic patients from a very small longitudinal cohort that were not merely
statistically similar to real patients but consistent with the underlying coagulation
biology. Our results provided evidence at each validation level that
SA-generated patients met this standard. Individual feature distributions were
reproduced with median MRE of 1.4\%, which MVN also achieved by construction. SA
additionally preserved the joint cross-visit covariance structure that MVN's
rank-deficient estimate could not represent, a difference made visible in the
eigenvalue spectrum where MVN introduced spurious variance in 194 null dimensions
while SA respected the low-rank geometry of the data. SA could amplify clinically
meaningful rare subgroups, generating 100 synthetic PCOS patients from only 3 real
ones, while preserving subgroup-specific signatures that MVN cannot capture. And an
independent ODE model of the coagulation cascade, calibrated
exclusively on real patients, confirmed that the factor-to-thrombin-generation mapping
in synthetic patients was biologically plausible, with 83--92\% cloud overlap and
Kolmogorov--Smirnov tests unable to distinguish the two populations ($p > 0.35$ for
all five TGA features).

The ability to generate validated synthetic cohorts from very small longitudinal
datasets has practical implications for maternal health research. Rare pregnancy
complications such as PCOS and preeclampsia are difficult
to study because assembling large, well-characterized longitudinal cohorts requires
years of recruitment across multiple clinical sites. Our results suggest that SA can
enable hypothesis generation, power analyses, and computational modeling studies for
these understudied populations by generating synthetic cohorts that preserve both the
statistical and mechanistic properties of the real data. The conditional generation
capability is particularly relevant. Rather than requiring researchers to collect
additional rare patients, SA can amplify existing small subgroups into cohorts large
enough for meaningful analysis.

Understanding why SA succeeds where parametric methods fail requires examining the
geometric structure of the problem. SA generates samples by interpolating between
stored patient profiles via the Hopfield
energy landscape operating in a PCA-reduced linear subspace, rather than fitting a
parametric distribution that must regularize or truncate high-dimensional structure. In
the $n < p$ regime that characterizes most small clinical studies, any parametric
distribution faces a fundamental identifiability problem. There are more parameters to
estimate than observations to estimate them from. SA sidesteps this by operating in a
PCA-reduced space where the memory-to-dimension ratio is favorable
($K/d_{\text{PCA}} \approx 1.28$) and by preserving the full directional structure of
the data through the attention mechanism. A natural concern is whether the $\beta^*$
selection procedure, originally developed for protein sequence
generation~\citep{Varner2026SAProtein}, transfers to clinical coagulation data. The entropy
inflection method identifies the phase transition in the Hopfield energy landscape,
which is a geometric property of $K$ patterns in $d$ dimensions, not a property of
what the features represent; for our dataset the empirical inflection
($\beta^* \approx 2.94$) was close to the theoretical prediction
($\sqrt{d} \approx 4.24$), and a sensitivity analysis confirmed that generation
quality degraded gracefully across $\beta/\beta^* \in [0.1, 3.0]$ (Table~\ref{stab:beta-sweep}). 
Similarly, varying the PCA variance retention threshold from 85\% to 99\%
($d_{\text{PCA}} = 13$ to 21) had minimal impact on generation quality. Median MRE
remained between 1.2\% and 1.8\% across all thresholds (Table~\ref{stab:pca-sensitivity}),
confirming that the 95\% threshold was not a critical design choice. The multiplicity
weighting for conditional generation introduced an additional consideration. Achieving
80\% attention on the PCOS subgroup required $\rho \approx 26.7$, reducing the
effective pattern count to $K_{\text{eff}} \approx 4.6$ (multiplicity parameters
for all three subgroups in Table~\ref{stab:multiplicity}),
which explains the larger MREs observed for PCOS features because the energy landscape
was effectively shaped by fewer than 5 independent patterns, pulling means toward the
population average. The direction-magnitude decomposition compounds this limitation.
For PCOS, magnitudes are drawn from only 3 empirical norms, providing minimal
diversity in the scale of generated samples. More broadly, the SA framework and its
multiplicity-weighted conditioning have been independently validated on discrete
protein sequence generation from small family
alignments~\citep{Varner2026SAProtein}, on steering generation toward functional
subsets such as binding peptides~\citep{Varner2026Multiplicity}, and on the
theoretical foundations connecting modern Hopfield networks to
attention-based generation~\citep{Alswaidan2026SA}. Across these domains, the
critical temperature prediction ($\beta^* \approx \sqrt{d}$) and the
multiplicity-weighted conditioning mechanism transfer without modification,
suggesting that the geometric properties exploited here are not specific to
clinical coagulation data but are general features of the Hopfield energy
landscape operating on small pattern sets.

An alternative precedent for validated longitudinal synthetic generation is
the VAMBN-MT framework of K\"uhnel et al.~\citep{Kuhnel2024}, which extends a
Variational Autoencoder Modular Bayesian Network with an LSTM time-encoder to
better capture cross-visit dependencies. They evaluated VAMBN-MT on the DONALD
nutritional cohort and on the Alzheimer's Disease Neuroimaging Initiative
(ADNI) cohort. Their study is instructive because it offers a regime contrast
to ours. With $N{=}1{,}312$ DONALD participants and 33 longitudinal variables,
their setting sits in the data-rich regime for which VAMBN-MT was designed,
where deep parametric models can be trained directly on the cohort. The
published configuration combines expert-curated module assignments,
Bayesian-network edge constraints, an LSTM-augmented HI-VAE per module, and
10{,}000 synthetic samples drawn from the trained model to reproduce a
published age-trend regression. In their evaluation, the best variant reported
a cross-visit correlation matrix relative Frobenius error of roughly 0.70, and
downstream time trends weaker than the strongest signals in the real data were
not reproduced at the sample sizes they tested, both honest indicators of how
hard this generative task is even in the data-rich regime. On the other hand,
our setting inverts both the number of patients and the number of features.
With $K{=}23$ patients across 216 features we operate two orders of magnitude
smaller and squarely in the $n < p$ regime, well below the cohort sizes for
which VAMBN-MT was developed. The non-parametric design
of SA addresses this regime in three ways that contrast with VAMBN-MT. First,
PCA-derived linear geometry replaces the modular HI-VAE architecture as the
mechanism for capturing cross-visit dependencies. Second, a generic
mechanistic-model validation pattern, transferable to any domain with a
calibrated ODE model, replaces the domain-specific dependency rules used in
their evaluation, such as graduation-status monotonicity, age-time arithmetic
identities, and macronutrient summation. Finally, multiplicity weighting
provides an inference-time conditional-generation lever that the VAMBN-MT
framework does not offer. These are complementary positions in the design
space rather than a head-to-head comparison. VAMBN-MT addresses the
moderate-cohort regime where deep generative models are feasible but
expert-specified modular structure helps capture cross-visit dependencies,
while SA addresses the small-cohort regime where parametric models cannot be
fit and the geometry of the few real patients must itself carry the generative
signal.

Beyond these statistical and geometric considerations, the mechanistic validation
introduced an approach that applies beyond coagulation. Rather than relying solely on statistical comparisons between real
and synthetic distributions, we tested whether synthetic patients produced biologically
consistent outputs when processed by an independent computational model, and went
further to show that a mechanistic model calibrated entirely on synthetic data
predicted real patient outcomes as well as one calibrated on real data (overall ratio
0.94$\times$ on held-out V2+V3 patients, with 11 random restarts per calibration).
This approach is available whenever a validated mechanistic model exists for the domain
of interest (pharmacokinetic models in drug development~\citep{MouldUpton2013},
physiological models in critical care~\citep{Chase2011}, tumor growth
models~\citep{Ribba2012}, or metabolic models~\citep{Yizhak2015,Shan2018}), and the requirement is not that the
model be a perfect predictor but that synthetic patients fall within the same
model-predicted envelope as real patients. The bootstrap Mann--Whitney analysis of
condition-specific features revealed that 5 of 24 (21\%) feature--condition pairs
were statistically distinguishable at equal sample sizes, but the mean relative errors
for these features remained small (5--15\%), and the detected differences likely
reflected SA's tendency to compress distributional tails rather than shift means. This
tail compression arises from two sources: the PCA dimensionality reduction, which
cannot fully represent higher-order moments of every feature, and the
direction-magnitude decomposition, which draws magnitudes from a finite empirical
distribution of $K{=}23$ norms. This is a trade-off. SA
sacrifices some tail fidelity in exchange for preserving the manifold geometry that
parametric methods cannot capture in the $n < p$ regime.

This study had several limitations. Our study used a single longitudinal
dataset of $K{=}23$ patients; further evaluation on independent datasets and
across different clinical domains is needed to establish generalizability. The PCA
dimensionality reduction assumed linear structure in the feature space; nonlinear
manifold learning methods may better capture complex clinical relationships, though at
the cost of reduced interpretability. The mechanistic validation under TF+TM
conditions revealed systematic ODE model bias, though the observation that real and
synthetic patients shared the same bias pattern was itself informative. The BZ2012
calibration required substantial departures from published rate constants, notably a
50-fold reduction in intrinsic tenase $k_{\text{cat}}$ and a 15-fold increase in
extrinsic tenase $k_{\text{cat}}$ (Table~\ref{stab:ode-calibration}). These deviations likely
reflect differences between the \textit{in vitro} conditions under which the
literature rate constants were measured (purified component systems) and the
plasma-based TGA used in this study, where phospholipid surfaces, factor
concentrations, and inhibitor profiles differ from the model
assumptions~\citep{Hockin2002,BrummelZiedins2012}. The purpose of the mechanistic
validation was not to obtain biochemically accurate rate constants but to test whether
the ODE model processed real and synthetic patients equivalently, a comparison that is
invariant to the absolute parameter values. The mechanistic validation was limited to
thrombin generation (BZ2012) and did not cover the fibrinolytic or viscoelastic
features included in the dataset; extending this validation to a mechanistic model of
fibrinolysis, currently under development, would provide additional evidence for these
feature categories. Finally, we did not validate synthetic patients against clinical
outcomes; our validation was restricted to statistical and mechanistic plausibility,
and prospective validation against real clinical decision-making remains an important
future direction.

\rone{A distinct limitation concerns privacy. Because the generator concentrates
probability mass near each of the $K{=}23$ stored patterns, synthetic patients lie
close to individual real records, and a nearest-neighbor membership-inference attack
separates cohort members from held-out patients with high accuracy. We find that this
behavior is intrinsic to the target distribution at this cohort size rather than an
artifact of the sampler or the operating temperature: an exact Metropolis-adjusted
sampler leaves it essentially unchanged, and lowering the inverse temperature improves
the distance to the closest real record only marginally, never to the level separating
real patients from one another, and at a measurable cost to cross-visit covariance
fidelity. We therefore do not present the method as privacy-preserving. It is not
differentially private, and the practical risk is bounded instead by the fully
de-identified nature of the source cohort and by intended uses, such as mechanistic
calibration and hypothesis generation, that do not require releasing patient-level
records. The burn-in and thinning used in the sampling diagnostic were derived once on
the full cohort and reused for the subset analyses, so the diagnostic characterizes the
operating regime rather than certifying each subset independently.}